% Reader question: is a constant audit rate the cheapest way to sustain
% deterrence over time, and by how much do the alternatives save? Claim:
% letting the audit rate fall as accumulated stake grows spends
% Theta(log T) instead of Theta(T); letting cheats surface on their own
% spends O(1), eventually needing zero new audits, but it is NOT uniformly
% cheaper -- at the chapter's own horizon T=200 it has spent more than the
% logarithmic schedule. Evidence: G=10, d=0.8, B=50, a=1, rho*=G/(dB)=0.25;
% flat spend at T=200 is 50.00; model A (loss-if-audited) is 25.58; model B
% (delayed revelation, rate r=0.05) is 35.08 at T=200, converging to
% aG^2/(2dBrv)=41.67 and reaching zero audit at t*=G/(rv)=333 (b2_tower.py,
% seed 20260816). Must distinguish: three different asymptotic classes
% (linear, logarithmic, bounded), and that "eventually free" (model B) is
% not the same claim as "cheaper at every horizon" -- model A is cheaper
% than model B at T=200 even though model B's total is bounded. Grammar: an
% xy plot, three schedules on one axis with the worked horizon marked,
% rejected: three numbers in a paragraph (states the totals but hides the
% shapes that explain them) and a bar chart at T=200 alone (hides that
% model B is still climbing past its rivals at that horizon).
% Five-second test: a reader says "B ends lowest eventually, but A is
% cheapest at the horizon this chapter actually uses."
\begin{figure}[H]
\centering
\pdfigurehue{pdviolet}%
\begin{tikzpicture}[pd figure]
\begin{axis}[pd axis,
  width=6.9cm,height=4.3cm,
  xmin=0,xmax=400,ymin=0,ymax=100,
  xtick={0,100,200,300,400},ytick={0,25,50,75,100},
  xlabel={period $T$},
  ylabel={cumulative audit spend}]
  \addplot[pd series,mark=none,domain=0:400,samples=2] {0.25*x};
  \addplot[pd focus series,mark=none,domain=0:400,samples=120]
    {(1*10/(0.8*0.6))*ln(1+0.6*x/50)};
  \addplot[pd ready series,mark=none,domain=0:333.33,samples=120]
    {(1/(0.8*50))*(10*x-(0.05*0.6/2)*x^2)};
  \addplot[pd ready series,mark=none,domain=333.33:400,samples=2] {41.67};
  \draw[pd guide] (axis cs:200,0) -- (axis cs:200,72);
  \node[pd label,anchor=south east,text width=2.3cm] at (axis cs:190,74) {$T{=}200$: this chapter's horizon};
  \node[pd datum] at (axis cs:200,50.00) {};
  \node[pd focus datum] at (axis cs:200,25.58) {};
  \node[pd ready datum] at (axis cs:200,35.08) {};
  \node[pd label,anchor=west] at (axis cs:405,90) {flat: $\Theta(T)$};
  \node[pd focus label,anchor=west] at (axis cs:405,60) {model A: $\Theta(\log T)$};
  \node[pd ready label,anchor=west] at (axis cs:405,41.67) {model B: $O(1)$};
  \draw[pd guide] (axis cs:333.33,0) -- (axis cs:333.33,10);
  \node[pd ready label,anchor=south west,text width=2.3cm] at (axis cs:337,2)
    {$t^\star{=}333$: audit rate hits zero};
\end{axis}
\end{tikzpicture}
{\linespread{1.18}\selectfont
\caption{Model A's audit rate falls, but its cumulative bill has no finite
ceiling. [verified, \protect\path{b2_tower.py}, seed 20260816]}
\label{fig:stp-audit-amortization}}
\end{figure}
